% Part: first-order logic

\documentclass[../../include/open-logic-part]{subfiles}

\begin{document}

\olpart{fol}{Logika Sepisanan}

\begin{editorial}
  Perangan iki nyinaoni sintaksis lan semantik logika sepisanan. Aturan
  pambentukan nemtokaké term, formula lan ukara. Struktur
  nentokaké interpretasi simbol, déné pangaji variabel lan
  relasi nyukupi nerangaké kapan formula utawa ukara bener.

  Kalkulus sekuèn, deduksi natural, tableau lan deduksi
  aksiomatis menehi cara kanggo mbuktèkaké akibat logis.
  Kasahihan lan kelengkapan nggandhèngaké derivasi karo
  akibat semantis. Teorema kekompakan lan L\"owenheim--Skolem
  mbatesi sipat kang bisa diandharaké ing basa iki.
\end{editorial}

\olimport[introduction]{introduction}

\olimport[syntax-and-semantics]{syntax}

\olimport[syntax-and-semantics]{semantics}

\olimport[models-theories]{models-theories}

\olimport[proof-systems]{proof-systems}

\iftag{prfSC}{%
  \olimport[sequent-calculus]{sequent-calculus}
}{}

\iftag{prfND}{%
  \olimport[natural-deduction]{natural-deduction}
}{}

\iftag{prfTab}{%
  \olimport[tableaux]{tableaux}
}{}

\iftag{prfAX}{%
  \olimport[axiomatic-deduction]{axiomatic-deduction}
}{}

\olimport[completeness]{completeness}

\olimport[beyond]{beyond}

\OLEndPartHook

\end{document}

